\scanpage{31}
\mspage{30}
l'image par \deleted{\unsure{(1.79)|(1.39)|(1.49)}} ($\unsure{*|\alpha}$)
\[
 \unsure{M|\mathcal M}\xrightarrow{\sim}\overset{\deleted{\unsure{\sim|\approx}}}{M}^{*}\subset\unsure{\mathcal P|\mathscr P|\mathfrak P|P}(M)
\]
d'une unique \inserted{partie} $F\subset M$, de sorte que $\Phi=\widetilde F$, i.e.
\[
 \deleted{\unsure{\mathrm{At}|\mathrm{AC}}}\qquad
 \Phi=\bigl\{\underset{\inserted{\substack{\shortparallel\\ M_{\ll x}}}}
 {\mathrm{Omb}(\deleted{\unsure{\widetilde X}}X)}\mid X\in F\bigr\}
\]
Ainsi, on peut aussi interpréter, grâce à $\ll$, $\unsure{\mathcal F|\mathfrak F|\mathscr F}_M$ en termes d'un\deleted{e} \deleted{partie} ens de parties \inserted{$\widetilde{\mathcal F}$} de $M$
\diagram{\begin{tikzcd}[ampersand replacement=\&,column sep=large,row sep=1pt]
 \mathcal F_M \arrow[r,"\sim"] \& \widetilde{\mathcal F}\\
 \cap \& \cap\\
 \mathcal P(\mathcal P(M) \& \mathcal P(M)
\end{tikzcd}}

% slanted
\margin{cf GF $\overline{\underline{\mathrm{IV}}}$ page 54 (Cor. 2) et page 55 (Remarque) Pour C3, cf page 58}
Ceci dit, \deleted{on demande que} \inserted{\unsure{Je veux|on veut|on veux} que $\widetilde{\mathcal F}\subset\mathcal P(M)$,}
\deleted{\unsure{soit associé à un}} \deleted{\unsure{satisfasse}} \inserted{\unsure{pour|muni|muni de}} \deleted{\unsure{cela|une|cette} relation d} la relation d'inclusion, satisfasse C1, C2, C3,
et \deleted{\unsure{identifiant|identifie} les \unsure{à}} pour $M$ \unsure{s'identifiant|s'identifie|s'identifier} à l'ens de ses objets \inserted{irréd.}

\deleted{\uline{\unsure{At|AE}}$^{*}$2 \unsure{bis}}\quad \inserted{On voudra d'abord que} L'ens. \unsure{des|de} parties $\widetilde{\mathcal F}$ \deleted{de} $\unsure{\subset}\mathcal P(M)$ soit \uline{préadmissible}, i.e. pour tout $X\in M$, il existe un plus petit élément de $\widetilde{\mathcal F}$ qui le contienne. En d'autres termes,

\uline{At$^{*}$2}$_{\unsure{0|o}}$\quad
\margin{\rule{0.45pt}{4\baselineskip}}
\inserted{Pour} tout $\deleted{\Phi\in\mathcal F_M}$ $M$-figure ensembliste $\Phi$ \deleted{\unsure{de|dans}} $\in\mathcal F_M$, et pour $A$ [$=\mathrm{Omb}(x)$!] \deleted{\unsure{de}} strate de $\Phi$ (i.e. $A\in\Phi$), il existe une \uline{plus petite figure} \inserted{$\Phi_A$} \deleted{\unsure{tous|dans} \unsure{des}} $\in\overset{\deleted{\sim}}{\mathcal F}_M$
% slanted
\margin{c'est la \uline{\unsure{multistrate}} définie par $A$}
(plus petite au sens de l'\uline{inclusion} \deleted{\unsure{)}} entre figures i.e. entre parties de $\mathcal P(M)$, ou \deleted{de} $M^{*}\subset\mathcal P(M)$) qui contienne $A$ comme élément (laquelle $\Phi_A$ sera donc $\subset\Phi$).

% slanted
\margin{C'est C1 pour $\unsure{\mathcal F_M|F_M}$}
\uline{At$^{*}$}3$_0$\quad Si $\Phi\in\mathcal F_M$, et si \deleted{\unsure{des}} $\Phi_i\in\mathcal F_M$, $\Phi_i\subset\Phi$, alors $\bigcup_{\unsure{i}}\Phi_i\in\mathcal F_M$ (stabilité par réunions \uline{conjointes} \inserted{dans $\mathcal F_M$})
